\documentclass[11pt,a4paper]{article}
\usepackage[margin=2.2cm]{geometry}
\usepackage{amsmath,amssymb,xcolor,enumitem,booktabs}
\usepackage[normalem]{ulem}
\usepackage[hidelinks]{hyperref}
\setlength{\parskip}{0.6em}\setlength{\parindent}{0pt}
\definecolor{newc}{HTML}{1f6f43}
\newcommand{\new}[1]{\textcolor{newc}{#1}}
\newcommand{\old}[1]{\textcolor{gray}{\sout{#1}}}
\newenvironment{greenblock}{\color{newc}}{}
\newcommand{\lref}[1]{{\small\texttt{#1}}}
\newcommand{\moved}[1]{{\small\textcolor{gray}{\emph{[moved from #1]}}}}
\title{Proposed rewrite, Appendix F, Demand drift}
\date{23 November 2026, version 2}
\begin{document}\sloppy
\maketitle

\section*{How to read this}
This proposal settles the open points of
\texttt{appendix\_F\_demand\_comments.md} (A1, A5 to A9, B1 to B3, C2 to C4,
D1 to D3, E1) and its three changes that would help the reader most. It
follows the decisions you have already made. The SKU group is a
\emph{category} ([A115]). The order gain stands beside the pick gain, from
the existing rows, with no rerun. The pilot, the initial run and the updated
run are told once each, in that order.

Part~1 is the whole appendix as it would read. Black text stays,
\new{green} is new and \old{struck} text goes. Text that only moves is black,
with a grey note saying where it came from. References are printed as their
labels, such as \lref{app:load:runs}. Every label of the appendix stays
(\lref{app:demand:*}, \lref{app:cal:drift}, \lref{fig:app:demandsweep}).
\lref{app:demand:sweep} now sits on the updated run, which is the run that
tests the values. One label is new, \lref{app:demand:first}, for the initial
run, and so is \lref{tab:app:demandgain}. Part~2 says where every number comes
from, and Part~3 holds the questions. Nothing below is in the thesis yet, and
Figure~F.1 is not redrawn.

Three things changed in the numbers, not only in the words.
\begin{itemize}[nosep]
  \item Every gain and cost is now the mean over seeds of the per-seed share,
        as \lref{eq:relgain} defines it. The old text used the ratio of the
        means for the initial run and for the put-away cost. So the initial
        run's pick gain at 32 agents reads 6 percent, not 5, its range 4 to
        12, not 3 to 13, and the put-away cost reads 9, 12, 13 and 14
        percent, not 10, 12, 13 and 15.
  \item The pilot now gives its numbers per starting layout. The old ``almost
        nothing'' and ``most of it back'' hold for 20 SKUs but not for every
        case (A6).
  \item The rows were written before [A96] and use $f_{\mathrm{up}} = 0.8$.
        They hold no served counts, so the keep-up check cannot be recounted.
        The lowest finished share of any row is 0.92 (Question~6).
\end{itemize}

\section*{Part 1. The appendix as it would read}

\subsection*{F.1 Summary \hfill\lref{app:demand:summary}}

Demand storage can only gain if demand gives it something to adapt to.
\old{The question is what demand that is, and whether the gain holds when
the demand settings change. Under uniform demand no placement beats another,
and under demand that never changes a placement computed once is hard to
beat. So demand in the experiments is skewed by a Zipf law over the SKUs, and
its ranking drifts by 5 percent of the SKUs every 200 steps. A pilot with
Token Passing showed that the drift is what made storage matter. The values
this sets in the thesis are the Zipf exponent of 1, the drift period of 200
steps and the drift step of 5 percent of the SKUs. The updated Token Passing
run shows that the gain holds over the whole range of each value at every
fleet size.}
\new{So the orders in the experiments follow demand that is skewed and
drifts. A few SKUs are ordered far more often than the rest, and which SKUs
these are changes over time. Two questions follow. Is the drift needed? And
how much does the gain of demand storage over fixed storage depend on the
values that set the demand?}

\begin{greenblock}
A pilot with Token Passing answers the first. A layout built once from
demand was hard to beat while demand held still. Once the ranking drifted,
that layout lost up to 15 percent of its throughput, and demand storage won
much of it back.

Two sweeps under Token Passing answer the second. Each varies the skew, the
drift period and the drift step one at a time. In the updated sweep demand
storage shortens an order by 14 to 18 percent against fixed storage at the
thesis values, and a pick by 10 to 14 percent. The gain holds at every swept
value and every fleet size.

The thesis uses a Zipf exponent of 1, a drift every 200 steps and a drift
step of 5 percent of the SKUs. These are the thesis's own values, taken from
the pilot. The sweeps test them and do not choose them.
\end{greenblock}

\subsection*{F.2 Setting \hfill\lref{app:demand:setting}}

An order asks for several SKUs, each on its own order line, and every line
becomes one pick or more (Section~\lref{sec:pf:tasks}, Orders and tasks).
\new{An order has 5 lines on average. Once all its lines have reached the
workstation, a shipment carries the order to a dock.}

\begin{greenblock}
The SKUs are grouped into categories of 20. Each category is served by four
workstations on neighbouring rack faces.
\end{greenblock}
Each seed splits the SKUs at random over the categories, and the split holds
for the whole run. \moved{the Conclusion}
\new{The warehouse grows with the fleet, by 2.5 SKUs and half a workstation
per agent (Section~\lref{sec:impl:instances}). So there are 1, 2, 4 and 8
categories at 8, 16, 32 and 64 agents.}

\new{A few SKUs are ordered far more often than the rest.} The SKUs are
ranked by popularity, and the SKU in rank $r$, counting from $0$, gets the
weight $1/(r+1)$, a Zipf law with exponent 1
(Section~\lref{sec:impl:instances}). A larger exponent, the skew, lets the
popular SKUs take a larger share of the orders. \moved{below}
\new{With another skew the weight is $1/(r+1)$ raised to the power of the
skew.}

\begin{greenblock}
Storage studies often skew demand by classes. Silva et al.\ (2020), for
example, let a small group of the SKUs take 80 percent of the orders, a middle
group 15 percent and the rest 5 percent. The Zipf law instead gives every rank
a weight of its own, so a single value, the skew, sets how uneven demand is.
The Zipf law is the thesis's own choice, and so is its exponent of 1, taken
from the pilot.
\end{greenblock}

An order first picks a workstation at random. It then draws its lines from
that workstation's category, each line a different SKU, by the weights of the
category's SKUs scaled to sum to one. So every category gets the same share
of orders, and a SKU's share of all lines is only roughly in proportion to
$1/(r+1)$. \old{A larger exponent, the skew, lets the popular SKUs take a
larger share of the orders.}

\begin{greenblock}
At 8 agents all 20 SKUs sit in one category, and all four workstations serve
it. An order then goes to any of the four, whatever SKUs it asks for. So
demand storage has no single place to aim a SKU at, and the best it can do is
a shelf close to all four on average. From 16 agents on, the lines of a SKU
only go to the four workstations of its category.
\end{greenblock}

Every 200 steps the ranking moves by 5 percent of the SKUs. Every SKU climbs
that many places, and the most popular ones drop to the bottom. With 20 SKUs
that is one place, and with 160 SKUs eight. This is the drift, with a period
of 200 steps and a step of 5 percent. The ranking then goes round once every
4000 steps at every fleet size.

No storage rule knows the demand law. A rule only sees the orders that have
already come in. Demand storage places arriving stock by the popularity it
estimates over the last storage period, and fixed storage sends each SKU back
to its home shelf (Section~\lref{sec:method:storage}). \new{Demand storage
aims a SKU at the four workstations of its category, which it learns from the
SKU's first order line.}

\old{The gain is how much demand storage shortens the mean service time of a
pick against fixed storage, paired by seed, as a share of fixed. It is
(\lref{eq:relgain}) with the pick in place of the order.}
\begin{greenblock}
Under one seed both rules receive the same orders, because the orders do not
depend on where the stock is. So the two rules are compared seed by seed. The
gain under a seed is how much demand storage shortens the mean service time
against fixed storage, as a share of fixed storage's time. The gain reported
is its mean over the seeds (\lref{eq:relgain}).

The gain is read for two times. The order time runs from an order's release
until its shipment reaches the dock. It is the result of the thesis
(\lref{eq:meanservice}). The pick time runs from a pick's release until its
units reach the workstation. Storage acts on it directly, because a pick
starts at a shelf.
\end{greenblock}

The demand model decides whether the main result can say anything about
storage. \old{Under uniform demand the gain is zero whatever the controller,
so a result would only show that the demand model gave storage nothing to
do.}
\new{Under uniform demand every SKU is ordered equally often, so demand
storage has no SKU to put ahead of another, and it is expected to gain little.
A result would then mostly show that the demand model gave storage little to
do. Uniform demand was not run.}

\subsection*{F.3 Method \hfill\lref{app:demand:method}}

\old{Two runs test the demand model. A pilot asks whether the drift is
needed, and a sweep asks how much the gain depends on the three values.}
\new{Three runs test the demand model, all with Token Passing. A pilot asks
whether the drift is needed. Two sweeps, an initial run and an updated run,
ask how much the gain depends on the three values.}
Token Passing is used because it needs no training, so the demand model is
judged before any controller learns on it. \moved{the sweep paragraph}

The pilot ran Token Passing on an earlier 8-agent warehouse of 94 vertices,
before the warehouses of Appendix~\lref{app:instances}, with a horizon of 2000
steps and 4 to 6 seeds\new{, at 0.6 tasks per step and with 20 and with 100
SKUs}. It compared fixed and demand storage with the ranking held fixed and
with the ranking moving \new{by 5 places} every 200 steps, and read the
throughput. \new{Five places is 25 percent of 20 SKUs and 5 percent of 100,
and the drift step of 5 percent comes from the 100 SKUs.} It ran under an
earlier storage model, in which the demand rule moved stock already on the
shelves, up to a cap per period.
\new{Each SKU count started from two layouts. The grouped layout put the
popular SKUs together on the shelves nearest the workstations, which is what
classic slotting builds once from demand (Heskett, 1963). The spread layout
spread them over all the shelves.}

The \old{sweep varies}\new{sweeps vary} one value at a time and \old{holds}\new{hold} the other two at the values of
the thesis. The skew takes 0.5, 1 and 1.5, the drift period 100, 200 and 400
steps, and the drift step 2.5, 5 and 10 percent of the SKUs. Each point runs
Token Passing under fixed and demand storage on the same seeds, each seed on
its own warehouse of the fleet's size, at the fleet's main load
(Appendix~\lref{app:load}). \old{Both rules see the same orders under a seed,
so the gain is paired by seed.} The put-away service time of demand storage
over that of fixed storage is read as the cost\new{, paired by seed in the
same way}, and the keep-up check of picks and put-aways as the check.
\old{Token Passing is used because it needs no training, so the demand model
is judged before any controller learns on it.} Both Token Passing runs used
20 seeds per value.

\begin{greenblock}
Each run lasts 10000 steps, and the first 5000 are not read. The measures
then start once the random starting stock has been replaced by the rules' own
pallets (Appendix~\lref{app:load:first}). At the thesis values the 5000 steps
that are read hold one and a quarter turns of the ranking. A drift period of
400 steps makes one turn 8000 steps, and so does a drift step of 2.5 percent
from 16 agents on. Those points see five eighths of a turn. The ranking still
moves every period, so demand storage still has to follow it, but not every
SKU reaches the top while the run is read.
\end{greenblock}

The sweep does not choose the values. A demand model tuned for the largest
gain would favour demand storage. The sweep shows instead how much the result
depends on the values.

\subsection*{F.4 Results \hfill\lref{app:demand:results}}

\old{The sweeps below come from two runs of Token Passing, an initial
validation run and an updated run after the model changes it led to. What
changed between them is in Appendix~\lref{app:load:runs}.}

\subsubsection*{F.4.1 Pilot: demand drift \hfill\lref{app:cal:drift}}

\old{In a pilot with Token Passing, the drift is what made storage matter.
With the ranking fixed, the demand rule gained almost nothing over fixed
storage. With the drift, fixed storage lost about 15\% of its throughput and
the demand rule won most of it back.}

\begin{greenblock}
A layout built once from demand is hard to beat while demand holds still.
From the grouped layout, demand storage gained 0.002 to 0.022 tasks per step
over fixed storage without the drift, 0 to 5 percent of a throughput near
0.43. From the spread layout it gained 0.040 to 0.045, about 10 percent,
because fixed storage kept that poorer layout and demand storage rearranged
it.

The drift puts the grouped layout out of date. With the drift, fixed storage
lost 15 percent of its throughput from the grouped layout and 8 to 10 percent
from the spread one. With 20 SKUs demand storage kept its throughput and
gained 0.078 over fixed storage from either layout, more than the whole loss.
With 100 SKUs it gained back 0.034 of the 0.069 lost from the grouped layout,
and 0.029 of the 0.033 lost from the spread one.
\end{greenblock}

\subsubsection*{\new{F.4.2 The initial run} \hfill\new{\lref{app:demand:first}}}

% model revision 12, sweep of 2026-10-27 (rows in git history of thesis-progress, 524ef5a)
The initial Token Passing run swept the demand settings at its own main loads
(Appendix~\lref{app:load:first})\old{, with 20 seeds per value. The drift
step was set as a share of the SKUs, 2.5, 5 and 10 percent}. At the thesis
values, demand storage shortens a pick by \old{9, 9, 5 and 10}\new{9, 9, 6
and 10} percent against fixed storage at 8, 16, 32 and 64 agents, and 16 to
19 of the 20 seeds agree.
\new{An order gains less, 7, 6, 2 and 4 percent. Every order then had a
single line, so its time was one pick and a shipment, and the shipment starts
at the workstation, not at a shelf.}

The gain holds at every value of every setting. It stays between \old{3 and
13}\new{4 and 12} percent at every fleet, and every seed keeps up. The skew
matters most. At a skew of 1.5 the few popular SKUs run out of stock. At 64
agents 28 percent of the orders then wait for stock, against 7 percent at a
skew of 1. The drift period and the drift step change the gain by a few
percent, with no pattern shared by the fleets.

Two settings did not do what the thesis says. With 30 and 50 SKUs, 5 percent
of the SKUs rounds to 2, so the ranking went round in 3000 and 5000 steps
instead of 4000. At 8 agents, with 20 SKUs, a step of 2.5 and of 5 percent
both rounded to one SKU and gave the same run.

In the initial run every order also had a single line, because the runs left
the number of lines per order unset. \moved{F.4.3}

\subsubsection*{F.4.3 Where an order \old{goes}\new{went} \hfill\lref{app:demand:workstation}}

% model revision 12, D29
The \old{same sweep}\new{initial run} showed that demand storage had no clear
place to aim for. Each order went to a workstation drawn at random, and the
workstations stand on rack faces all over the warehouse. Demand storage still
put a popular SKU next to the nearest workstation, but the order then went to
any of them. Over five warehouses per fleet, the shelves the rule ranked best
were about 3 steps from their nearest workstation. From the workstation an
order actually used they were 11, 16, 24 and 34 steps away, at 8, 16, 32 and
64 agents. An average shelf was 12.5, 17, 25 and 35 steps away. So the best
shelves were only 3 to 10 percent closer in practice.

\old{In the initial run every order also had a single line, because the runs
left the number of lines per order unset.}

\subsubsection*{\new{F.4.4 What changed before the updated run}}

\new{The initial run showed two faults in the demand and one setting that was
missing. Each led to a change. The full list of changes between the runs is
in Appendix~\lref{app:load:runs}.} \moved{the Conclusion}
\old{Three things change, each for a reason the sweep showed.}
\begin{itemize}
  \item The number of SKUs is a multiple of 20. Then 5 percent of the SKUs is
  a whole number, and the ranking goes round in exactly 4000 steps at every
  fleet size, as the thesis says.
  \item Each SKU belongs to a category of 20 SKUs, served by four
  workstations on neighbouring rack faces, and an order takes all its lines
  from the category of its workstation. \old{A SKU then has one place where
  it is needed, so storing it close to that place means something.}
  \new{From 16 agents on, a SKU's lines then go to one group of four
  workstations, so storing it close to them means something. Demand storage
  now scores a shelf by its mean distance to the four workstations of the
  SKU's category, not by the nearest workstation.} \old{Each seed splits the
  SKUs at random over the categories, and the split holds for the whole run.
  An order picks its workstation at random.}
  \item \new{The missing setting.} Orders have 5 lines on average, as the
  thesis says, instead of the single line of the \old{sweep}\new{initial
  run}.
\end{itemize}
\new{The stock that ran out at a skew of 1.5 led to a larger safety stock,
the one of the most popular rank for every SKU
(Appendix~\lref{app:cal:safety}).}

\subsubsection*{F.4.5 \old{Demand skew and drift}\new{The updated run} \hfill\lref{app:demand:sweep}}

% step 12 second sweep, 2026-10-31 (thesis-progress results/F_demand/README.md)
The updated Token Passing run repeated the sweep with the categories of
Appendix~\lref{app:load:next}, at the new main loads. At the thesis values
demand storage now shortens \new{an order by 14, 18, 18 and 15 percent and}
a pick by 11, 14, 13 and 10 percent at 8, 16, 32 and 64 agents\old{, and the
gain stays between 9 and 15 percent over the whole range}\new{, as
Table~F.1 shows. Over the whole range the order gain stays between 12 and 22
percent and the pick gain between 9 and 15}. Every seed keeps up. \new{These
rows were counted before the served share of Appendix~E. The earlier count
took only the work released in the measured steps, so it is never above the
served share, apart from the few lines that wait for stock (at most 0.35
percent). The lowest share of any seed is 0.92, so every seed also keeps up at
the threshold of 0.9.}

\begin{greenblock}

\begin{center}
\small
Table F.1: Demand against fixed storage at the thesis values, updated run.
Mean over 20 paired seeds, with its standard error.
\hfill\lref{tab:app:demandgain}\\[0.4em]
\begin{tabular}{rrrrrrr}
\toprule
$m$ & SKUs & categories & order gain (\%) & pick gain (\%) & put-away cost (\%) & travel per pick \\
\midrule
8  & 20  & 1 & $14.1 \pm 2.8$ & $10.9 \pm 2.7$ & $9.1 \pm 2.6$  & 17.0 to 13.2 \\
16 & 40  & 2 & $17.5 \pm 2.9$ & $13.9 \pm 2.6$ & $11.6 \pm 1.7$ & 19.5 to 15.0 \\
32 & 80  & 4 & $17.7 \pm 1.4$ & $13.2 \pm 1.5$ & $12.5 \pm 1.5$ & 22.2 to 17.3 \\
64 & 160 & 8 & $14.6 \pm 1.3$ & $10.2 \pm 1.3$ & $14.3 \pm 1.0$ & 25.1 to 20.3 \\
\bottomrule
\end{tabular}
\end{center}
\end{greenblock}

\new{Demand storage cuts the travel of a pick by 19 to 23 percent at the
thesis values,} for example from 17 to 13 steps at 8 agents. In the initial
run it cut the \old{driving}\new{travel} by 9 to 11 percent. The service time
gains less\old{, 10 to 14 percent,} because under Token Passing a pick mostly
waits for its workstation and \old{driving}\new{travel} is a small part of
its time (Appendix~\lref{app:load:next}). \moved{F.4.3}
\old{In the updated run every SKU belongs to a category served by four
neighbouring workstations, so storage has a place to aim for.}

\begin{greenblock}
The categories cannot explain the larger travel cut at 8 agents. There all
SKUs share one category, and an order may still go to any workstation. Yet
the travel cut is 22 percent, as large as at 16 and 32 agents. Other things
changed between the runs at the same time, the score by the mean distance to
the category's workstations, 5 lines per order and pallets of 10 units. The
runs do not separate them. So they show that storage gains more after the
changes, not how much of it the categories give.

At 8 agents the gain is among the smallest, with 64 agents. The travel cut at
8 agents is as large as at the larger fleets, so the single category does not
show in the travel, and these runs do not say why the gain there is smaller.
\end{greenblock}

The cost is in the put-aways. Demand storage sends stock to the shelves near
the workstations, so its put-aways take \old{10, 12, 13 and 15}\new{9, 12, 13
and 14} percent longer than those of fixed storage at the thesis values, and
\old{9 to 16}\new{7 to 16} percent longer over the whole range. The cost
grows with the fleet.

The skew no longer matters much. It moves the \new{pick} gain by at most 3.5
points\new{ and the order gain by at most 3}, because the stock no longer
runs out. \old{Between 0.1 and 0.3 percent of the lines wait for stock at
every value.}\new{Fewer than 0.4 percent of the lines wait for stock at any
value.} The drift period and the drift step again change the gain by only a
few points\new{, at most 3 for a pick and 5 for an order, with no pattern
shared by the fleets}. At 8 agents a step of 2.5 and of 5 percent still round
to the same single SKU.

% Figure: thesis-progress/scripts/images/appendix_load.py, demand_sweep()
Figure F.1 \lref{fig:app:demandsweep}, caption: \old{The gain of demand over
fixed storage, its put-away cost and the keep-up check as the demand model
changes, under Token Passing.}
\new{The order and pick gain of demand over fixed storage and its put-away
cost as the demand model changes, under Token Passing.} (22 words. The
figure is redrawn, order gain, pick gain and put-away cost in percent, no
keep-up row.)

\subsection*{F.5 Conclusion \hfill\lref{app:demand:conclusion}}

\old{The need for drift is settled. Without it the demand rule gained almost
nothing over fixed storage, so the experiments use demand that is skewed and
drifts.}
\new{The drift is needed. A layout built once from demand is hard to beat
while demand holds still, and the drift puts it out of date. So the
experiments use demand that is skewed and drifts.}

\old{The values themselves, an exponent of 1 and a drift every 200 steps,
have no source. They are the values of the pilot, and the sweep shows how
much the gain depends on them rather than choosing them.}
\begin{greenblock}
The three values are the thesis's own. Skewed demand is usual in storage
studies (Silva et al., 2020), but none of the sources the thesis uses gives a
Zipf exponent or a drift for this setting. The exponent of 1 and the drift
every 200 steps are the values of the pilot. The drift step of 5 percent is
the pilot's step of 5 places with 100 SKUs. The drift period is tied to the
storage period of 100 steps, which is half of it, so demand storage sees every
move of the ranking within one period (Appendix~\lref{app:cal:period}).
\end{greenblock}

\old{The initial run shows that the gain holds over the whole range at every
fleet size, so the exponent of 1, a drift period of 200 steps and a drift
step of 5 percent are kept.}

\old{Three things change, each for a reason the sweep showed.} \moved{here to F.4.4, list not repeated}

The \new{two sweeps keep the values.} \old{updated run, with these changes,
confirms the values.} The gain holds over the whole range at every fleet
size\old{, between 9 and 15 percent,}\new{. In the updated run an order gains
12 to 22 percent and a pick 9 to 15 percent,} and the skew no longer matters
much because the stock no longer runs out.

\new{At 8 agents all SKUs share one category, so demand storage has no single
place to aim for. Its travel cut is still as large as at the larger fleets.}

\section*{Part 2. Where the numbers come from}
All gains and costs are the mean over seeds of the per-seed share, as
\lref{eq:relgain}. Fields \texttt{mean\_order\_completion} (order),
\texttt{mean\_service\_time} (pick), \texttt{mean\_putaway\_service\_time},
\texttt{mean\_travel\_time}, \texttt{num\_lines\_waited}.
\begin{itemize}[nosep]
  \item Updated run. \texttt{thesis-progress/results/F\_demand/fleet\_<m>/demand.*/load\_sweep.jsonl}.
        At the thesis values the order gain is 14.1, 17.5, 17.7, 14.6 and the pick gain
        10.9, 13.9, 13.2, 10.2 (matching the comments file and the README). Seeds that
        agree, order 18, 19, 20, 20, pick 17, 19, 20, 20. Over every swept value, order
        12.5 to 21.9, pick 8.8 to 14.5, put-away cost 7.0 to 15.9, travel cut 17.4 to
        23.6. Lines waiting for stock 0.02 to 0.35 percent. Lowest finished share 0.92.
  \item Initial run. The same files at \texttt{524ef5a}. Pick 9.0, 9.3, 5.8, 9.7 (the
        ratio of means gives 9.2, 9.5, 5.5, 9.8, the old text). Order 6.8, 5.9, 2.4, 3.6.
        Travel cut 11.0, 10.5, 9.0, 9.1. Pick range 3.8 to 12.3. At 64 agents 28.4 and
        7.0 lines waited per 100 orders at skews 1.5 and 1.
  \item Pilot. \texttt{decisions/done/demand\_probe\_2026-10-17.md} and its rows
        \texttt{quick/summary.md}, restored to
        \texttt{thesis-progress/results/F\_demand/pilot\_2026-10-17/}.
  \item Categories and the order draw. \texttt{multi\_agent\_env.py},
        \texttt{\_draw\_order\_lines}, \texttt{popularity\_ranking},
        \texttt{zipf\_weights}. The drift step is
        \texttt{max(1, round(share * SKUs))} in \texttt{experiment\_config.py}.
  \item Skewed demand by classes. Silva et al.\ 2020
        (\texttt{silva2020integrating\_slap\_opp}), ``Random, Skewed 50/30/20,
        Skewed 80/15/5''. No paper in \texttt{papers/} uses a Zipf law.
\end{itemize}

\section*{Part 3. Answered by the author (23 November)}
\begin{enumerate}
  \item Figure F.1 is redrawn as proposed.
  \item Uniform demand is not run. The text says what is expected and that
        it was not run.
  \item The sentence on why an order gains more than a pick is dropped. It
        was not measured, and an order's lines are separate tasks that
        different agents can serve at the same time.
  \item Keep-up is fixed by the argument above, without a rerun.
  \item The pilot rows are restored from the Trash.
\end{enumerate}

\section*{Part 4. Question left}
\begin{enumerate}
  \item \textbf{The 8-agent argument.} The rows do not show the single
        category hurting at 8 agents. The travel cut there is 21.8 percent,
        as large as at 16 (22.7) and 32 (22.0), and the gain is just as small
        at 64 agents. The draft says the runs do not explain the smaller
        gain at 8 agents. Agree?
        \\ \hspace*{1em}$\square$ agree \quad $\square$ other \quad Note:
\end{enumerate}

\end{document}
